\documentclass[10pt,a4paper]{amsart}
\usepackage{amsmath,amssymb,amsthm,graphicx}
\usepackage{verbatim}
\newenvironment{reference}{\comment}{\endcomment}
%\newenvironment{reference}{}{}
\newenvironment{slogan}{\comment}{\endcomment}
\newenvironment{history}{\comment}{\endcomment}

% For commutative diagrams we use Xy-pic
\usepackage[all]{xy}

% We use 2cell for 2-commutative diagrams.
\xyoption{2cell}
% Theorem environments.
%
\theoremstyle{plain}
\newtheorem{theorem}[subsection]{Theorem}
\newtheorem{proposition}[subsection]{Proposition}
\newtheorem{lemma}[subsection]{Lemma}

\theoremstyle{definition}
\newtheorem{definition}[subsection]{Definition}
\newtheorem{example}[subsection]{Example}
\newtheorem{exercise}[subsection]{Exercise}
\newtheorem{situation}[subsection]{Situation}

\theoremstyle{remark}
\newtheorem{remark}[subsection]{Remark}
\newtheorem{remarks}[subsection]{Remarks}

\numberwithin{equation}{subsection}

% Macros
%
\def\lim{\mathop{\mathrm{lim}}\nolimits}
\def\colim{\mathop{\mathrm{colim}}\nolimits}
\def\Spec{\mathop{\mathrm{Spec}}}
\def\Hom{\mathop{\mathrm{Hom}}\nolimits}
\def\Ext{\mathop{\mathrm{Ext}}\nolimits}
\def\SheafHom{\mathop{\mathcal{H}\!\mathit{om}}\nolimits}
\def\SheafExt{\mathop{\mathcal{E}\!\mathit{xt}}\nolimits}
\def\Sch{\mathit{Sch}}
\def\Mor{\mathop{\mathrm{Mor}}\nolimits}
\def\Ob{\mathop{\mathrm{Ob}}\nolimits}
\def\Sh{\mathop{\mathit{Sh}}\nolimits}
\def\NL{\mathop{N\!L}\nolimits}
\def\CH{\mathop{\mathrm{CH}}\nolimits}
\def\proetale{{pro\text{-}\acute{e}tale}}
\def\etale{{\acute{e}tale}}
\def\QCoh{\mathit{QCoh}}
\def\Ker{\mathop{\mathrm{Ker}}}
\def\Im{\mathop{\mathrm{Im}}}
\def\Coker{\mathop{\mathrm{Coker}}}
\def\Coim{\mathop{\mathrm{Coim}}}

% Boxtimes
%
\DeclareMathSymbol{\boxtimes}{\mathbin}{AMSa}{"02}

%
% Macros for moduli stacks/spaces
%
\def\QCohstack{\mathcal{QC}\!\mathit{oh}}
\def\Cohstack{\mathcal{C}\!\mathit{oh}}
\def\Spacesstack{\mathcal{S}\!\mathit{paces}}
\def\Quotfunctor{\mathrm{Quot}}
\def\Hilbfunctor{\mathrm{Hilb}}
\def\Curvesstack{\mathcal{C}\!\mathit{urves}}
\def\Polarizedstack{\mathcal{P}\!\mathit{olarized}}
\def\Complexesstack{\mathcal{C}\!\mathit{omplexes}}
% \Pic is the operator that assigns to X its picard group, usage \Pic(X)
% \Picardstack_{X/B} denotes the Picard stack of X over B
% \Picardfunctor_{X/B} denotes the Picard functor of X over B
\def\Pic{\mathop{\mathrm{Pic}}\nolimits}
\def\Picardstack{\mathcal{P}\!\mathit{ic}}
\def\Picardfunctor{\mathrm{Pic}}
\def\Deformationcategory{\mathcal{D}\!\mathit{ef}}
\usepackage{cleveref}
\crefname{originalremark}{original remark}{original remarks}
\Crefname{originalremark}{Original remark}{Original remarks}
\makeatletter
\newlabel{remark-idempotent-symmetry}{{0H9N}{source}{}{sourcepointer-1}{}}
\newlabel{remark-idempotent-symmetry@cref}{{[originalremark][0][]0H9N}{[1][0][]source}{}{}{}}
\makeatother
\crefname{originallemma}{original lemma}{original lemmas}
\Crefname{originallemma}{Original lemma}{Original lemmas}
\makeatletter
\newlabel{lemma-kernel-mono}{{0E43}{source}{}{sourcepointer-2}{}}
\newlabel{lemma-kernel-mono@cref}{{[originallemma][0][]0E43}{[1][0][]source}{}{}{}}
\makeatother
\title{The kernel and the cokernel of an idempotent endomorphism}
\author{The Stacks Project Authors}
\date{stacks-project:homology.tex:local-snapshot-20261012; GNU FDL 1.2 or later}
\begin{document}
\maketitle
\noindent\textit{Editorial note.} All delivered mathematical text is the author's own, copied from the pinned Stacks Project chapter file homology.tex; no statement or proof step was rewritten, reordered or removed.\par\medskip
\noindent\textit{Editorial note.} The proof is the author's complete argument for the first assertion, including both halves of the universal property of the cokernel; the closing sentence 'By duality, symmetry ..., (2), (3), and (4) follow as did (1).' is the author's own reduction of the remaining three assertions to the one written out in full inside the counted span, and it is preserved verbatim.\par\medskip
\noindent\textit{Editorial note.} Every labelled reference occurring inside the counted spans is carried in assembly.external\_references as a non-load-bearing navigation pointer (same-chapter ones at their definition site, cross-chapter ones at the invocation inside the counted span), so no load-bearing mathematics is replaced by a stub. No bibliography citation occurs inside the counted spans.\par\medskip
\noindent\textit{Editorial note.} The same-name check for this card scanned every staged Stacks chapter of this repository -- all sources/stacks-*/pinned-original-source.tex files, the card's own chapter excluded -- and the labels found are recorded in dedup.same\_name\_cross\_chapter on the card.\par\medskip
\noindent\textit{Editorial note.} The preamble is the chapter's own native preamble excerpt already used for this source: the comment-type reference, slogan and history environments (native preamble lines 8-13), the Xy-pic setup (native preamble lines 14-18) and the native theorem declarations with the mathematical macros (native preamble lines 154-217, the actual last line of the file).\par\medskip
\noindent\textit{Allowed standard prerequisites (editorial).}\begin{itemize}
\item The complementary idempotent. Same chapter, Remark remark-idempotent-symmetry. Supplies that for an idempotent f the endomorphism g = id - f satisfies f g = g f = 0 and g g = g, so the statement may be applied with the roles of f and g exchanged.
\item A kernel map is a monomorphism. Same chapter, Lemma lemma-kernel-mono. Supplies that the kernel map i can be cancelled on the left, which turns the identity i p i = i into p i = id.
\end{itemize}\medskip
\ifdefined\hypertarget\hypertarget{sourcepointer-1}{}\fi
\noindent\textit{Source reference.} Original remark 0H9N: sources/stacks-homology/pinned-original-source.tex, source lines 575--575. Navigation pointer invoked at proof line 589; the counted span contains the pointer only and the cited result is not reproduced inside it, so the pointer is registered as a non-load-bearing navigation stub in the established form for Stacks material. No bibliography citation is involved.\par\medskip
\ifdefined\hypertarget\hypertarget{sourcepointer-2}{}\fi
\noindent\textit{Source reference.} Original lemma 0E43: sources/stacks-homology/pinned-original-source.tex, source lines 492--492. Navigation pointer invoked at proof line 530; the counted span contains the pointer only and the cited result is not reproduced inside it, so the pointer is registered as a non-load-bearing navigation stub in the established form for Stacks material. No bibliography citation is involved.\par\medskip
\begin{lemma}
\label{lemma-idempotent-kernel-cokernel}
Let $\mathcal{A}$ be a preadditive category,
$x$ be an object of $\mathcal{A}$,
and $f : x \to x$ be idempotent.
Write $g = \text{id} _ x - f$
and recall by Remark \ref{remark-idempotent-symmetry}
that $f \circ g = g \circ f = 0$.
\begin{enumerate}
\item If $f$ has kernel $i : \Ker (f) \to x$
and $p$ is the unique morphism such that $g = i \circ p$,
then $p \circ i = \text{id} _ {\Ker (f)}$
and $p$ is the cokernel of $f$.
\item If $f$ has cokernel $p : x \to \Coker (f)$
and $i$ is the unique morphism such that $g = i \circ p$,
then $p \circ i = \text{id} _ {\Coker (f)}$
and $i$ is the kernel of $f$.
\item If $g$ has kernel $j : \Ker (g) \to x$
and $q$ is the unique morphism such that $f = j \circ q$,
then $q \circ j = \text{id} _ {\Ker (g)}$
and $q$ is the cokernel of $g$.
\item If $g$ has cokernel $q : x \to \Coker (g)$
and $j$ is the unique morphism such that $f = j \circ q$,
then $q \circ j = \text{id} _ {\Coker (g)}$
and $j$ is the kernel of $g$.
\end{enumerate}
\end{lemma}

\begin{proof}
As for (1), compute first that $i = (f + g) \circ i = g \circ i$.
By the construction of $p$, this implies that
$i \circ p \circ i = g \circ i = i \circ \text{id} _ {\Ker (f)}$.
By Lemma \ref{lemma-kernel-mono} we conclude that
$p \circ i = \text{id} _ {\Ker (f)}$.

\medskip\noindent
It remains to show that $p : x \to \Ker (f)$
satisfies the universal property of the cokernel of $f$,
i.e., that given a morphism $a : x \to y$ satisfying $a \circ f = 0$
there exists a unique morphism $b : \Ker (f) \to y$
such that $a = b \circ p$.
Existence follows by first computing that $a = a \circ (f + g) = a \circ g$
and concluding by the construction of $p$
that $a = a \circ g = a \circ i \circ p$
(so that we can take $b = a \circ i$).
Uniqueness follows from that $b = b \circ p \circ i = a \circ i$
(so that $a$ determines $b$).

\medskip\noindent
By duality,
symmetry (i.e., that $g$ is itself idempotent;
cf. Remark \ref{remark-idempotent-symmetry}),
and duality and symmetry respectively, (2), (3), and (4) follow as did (1).
\end{proof}

\end{document}
